\begin{solucion}
Sea $L/K$ una extensión de grado $2$. En primer lugar, como $[L:K]$ es finito, $L/K$ es una \emph{extensión algebraica} de $K$ (es decir, cada elemento de $L$ es raíz de algún polinomio no nulo con coeficientes en $K$). En particular, existe al menos un elemento $\alpha \in L$ que no pertenece a $K$ (pues $L \neq K$) y $\alpha$ es algebraico sobre $K$. Denotemos por $m_{\alpha,K}(x)$ el \emph{polinomio minimal} de $\alpha$ sobre $K$, definido como el único polinomio monico de menor grado en $K[x]$ que tiene a $\alpha$ como raíz. 

Ahora, consideremos la cadena de extensiones $K \subseteq K(\alpha) \subseteq L$. Por el \emph{teorema de la torre} (multiplicación de grados), se cumple 
\[ [L:K] = [L:K(\alpha)] \,[K(\alpha):K] . \] 
Sabemos además que $[L:K] = 2$. Dado que $\alpha \notin K$, el grado $[K(\alpha):K]$ es mayor que $1$, por lo que forzosamente $[K(\alpha):K] = 2$. Esto implica que $[L:K(\alpha)] = 1$, es decir, $L = K(\alpha)$. En otras palabras, $L$ es una \emph{extensión simple} generada por el elemento $\alpha$. Además, por teoría de extensiones simples, el grado de $K(\alpha)$ sobre $K$ coincide con el grado de $m_{\alpha,K}(x)$. Por lo tanto $\deg(m_{\alpha,K}) = [K(\alpha):K] = 2$. Denotemos $f(x) := m_{\alpha,K}(x)$; entonces $f(x) \in K[x]$ es un polinomio irreducible de grado $2$. 

Sea $f(x) = x^2 + bx + c$, con $b,c \in K$. Como $f$ es irreducible sobre $K$, sus raíces $\alpha$ y $\beta$ no están en $K$ (al menos una de ellas, $\alpha$, está en $L$ por construcción). Observemos que la suma de las raíces vale $\alpha + \beta = -b$, que pertenece a $K$. En consecuencia, 
\[ \beta = -\,b - \alpha. \] 
Puesto que $-b \in K$ y $\alpha \in L$, se deduce que $\beta \in K(\alpha) = L$. Es decir, \emph{ambas} raíces $\alpha$ y $\beta$ de $f(x)$ pertenecen a $L$. (En el caso degenerado de que $f(x)$ tuviera una única raíz doble $\alpha$ --lo cual sólo ocurre si la característica de $K$ es $2$ y $f(x)=(x-\alpha)^2$--, esa raíz $\alpha$ también pertenece a $L$ por su definición. En cualquier caso, $L$ contiene todas las raíces de $f$.) Por lo tanto, $f(x)$ se factoriza completamente como producto de polinomios lineales en $L[x]$: 
\[ f(x) \;=\; (x-\alpha)\,(x-\beta) , \] 
con $\alpha,\beta \in L$. 

Hemos demostrado entonces que $L$ contiene a todas las raíces de $f(x)$ y, de hecho, $L = K(\alpha,\beta)$. Más aún, $L$ es la menor extensión de $K$ que contiene dichas raíces, ya que cualquier otro subcuerpo de $L$ que contenga a $\alpha$ (o a $\beta$) debe contener también a $\beta$ (o a $\alpha$) porque $\beta = -b-\alpha$, y por ende contendría a $K(\alpha)=L$. En conclusión, $L$ es el \emph{cuerpo de descomposición} de $f(x)$ sobre $K$, como queríamos demostrar.

\begin{center}
\begin{tabular}{|l|l|}
\hline 
\textbf{Concepto clave} & \textbf{Referencia (Archivo PDF, página)} \\ \hline 
Extensión simple & Clase 5.pdf, pág. 1 \\ \hline
Polinomio minimal & Clase 5.pdf, pág. 2 \\ \hline
Grado de extensión finita & Clase 6.pdf, pág. 1 \\ \hline
Teorema de la torre & Clase 6.pdf, pág. 2 \\ \hline 
Extensión finita $\implies$ algebraica & Clase 7.pdf, pág. 1 \\ \hline 
Cuerpo de descomposición & Clase 11.pdf, pág. 1 \\ \hline 
\end{tabular}
\end{center}
\end{solucion}